% Resumen (150 a 250 palabras, tres o cuatro cifras como mucho)
\section*{Resumen}
\addcontentsline{toc}{section}{Resumen}

Este informe reúne el contexto regional del problema que trabaja DespachoX, pensado para quien no vive en Manizales. El 10 de agosto de 2026, un sismo de magnitud 7,4 con epicentro en el Chocó llevó al Gobierno a declarar la emergencia en Caldas y en otros departamentos. En Manizales, el censo que cerró el 24 de agosto registró 8.240 personas damnificadas.

La comida llegó a la ciudad. La red de bancos de alimentos usó a Manizales como punto de acopio y, sumando todo su corredor, movió más de 2.000 toneladas en la primera semana. La dificultad apareció al repartir. Dos días después seguían cerradas la vía Pereira–Chinchiná y la vía a Fresno por el Alto de Letras. Los cortes de energía dañaron la cadena de frío, y la red señaló la llegada a las zonas rurales como su mayor reto.

DespachoX propone una mesa de despacho que recibe por chat los pedidos de los barrios, recalcula rutas cuando un conductor reporta un cierre y sigue cada paquete hasta su entrega. Una persona aprueba cada plan. En simulación, el sistema supera con claridad al despacho por orden de llegada en entregas a tiempo a puntos críticos. Frente a un despachador con oficio empata en ese indicador, así que el equipo sitúa su aporte en poner esa pericia al alcance de cualquier voluntario.

Ninguna cifra ha sido verificada todavía por una persona del equipo. La pregunta abierta que más pesa es cómo despacha hoy el Banco de Alimentos de Manizales.

\vspace{1.2em}
{\small\tableofcontents}
\clearpage
